%-*-latex-*-
%-*-latex-*-
\newcommand\COURSE{ciss245}
\newcommand\ASSESSMENT{q6701}
\newcommand\ASSESSMENTTYPE{Quiz}
\newcommand\POINTS{\textwhite{xxx/xxx}}

\input{myquizpreamble}
\input{yliow}
\input{\COURSE}
%-*-latex-*-

%-*-latex-*-

%-*-latex-*-
\renewcommand\TITLE{\ASSESSMENTTYPE \ \ASSESSMENT}

\renewcommand\EMAIL{}
\renewcommand\AUTHOR{eagraham1@cougars.ccis.edu}

\textwidth=6in
\begin{document}
\topmatterthree

%------------------------------------------------------------------------------
\nextq
Add code to your Rational class so that the following works:
\begin{Verbatim}[frame=single, fontsize=\footnotesize]
Rational f(5, 2); // f models 5/2
int i = int(f);   // i is 2
int j = f;        // THIS WILL FAIL
\end{Verbatim}
\ANSWER
\begin{answercode}
class Rational
{
public:
    Rational(int n, int d)
        : n_(n), d_(d)
    {}

    operator int() const
    {
      return n_ / d_;
    }

    // other methods not shown
private:
    int n_, d_;
}
\end{answercode}

%------------------------------------------------------------------------------
\nextq
Correct the code below so that it produces the following output:
\begin{Verbatim}[frame=single,fontsize=\footnotesize]
<Werewolf john>
<Lucan john>
\end{Verbatim}
\begin{answercode}
#include <iostream>
#include <string>

class Lucan
{
public:
    Lucan(const std::string name)
        : name_(name)
    {}

    std::string name_;
};

std::ostream & operator<<(std::ostream & cout, const Lucan & l)
{
    cout << "<Lucan " << l.name_ << '>';
    return cout;
}

class Werewolf
{
public:
    Werewolf(const std::string name)
        : name_(name)
    {}

    operator Lucan() const
    {
        return Lucan(name_);
    }

    std::string name_;
};

std::ostream & operator<<(std::ostream & cout, const Werewolf & w)
{
    cout << "<Werewolf " << w.name_ << '>';
    return cout;
}

int main()
{
    Werewolf w("john");
    std::cout << w << '\n';
    std::cout << Lucan(w) << '\n';
    return 0;
}
\end{answercode}

%------------------------------------------------------------------------------
%-*-latex-*-
\newpage
\input{instructions}
\end{document}

